% basic part begin
\documentclass[12pt,a4paper]{extarticle}

\usepackage{fontspec}
\usepackage[english]{babel}
\usepackage{amsmath,amsfonts,amssymb}
\usepackage[dvipsnames]{xcolor}
\usepackage{enumitem}
\usepackage{csquotes}
\usepackage{microtype}
\usepackage{fontawesome}
\usepackage{hyperref}
\usepackage{enumitem}

\usepackage[utf8]{inputenc}
\usepackage[OT1]{fontenc}
\usepackage[misc]{ifsym}
\usepackage{makeidx}
\usepackage{wasysym}
\usepackage{vwcol}
\usepackage{tgbonum}
\usepackage{tabularx}
\usepackage{array}
\usepackage{ragged2e}
\usepackage{changepage}

\hypersetup{
    colorlinks=false,
    pdfborder={0 0 0}
}

% basic part end
\usepackage{multicol}
\usepackage[left=1cm,right=1cm,top=1cm,bottom=1cm]{geometry}

\newcommand{\cvsection}[1]{%
    \vspace{1mm}
    \par\noindent\hbox{\Large{\textcolor{Plum}
    {\textbf{#1}}}}\vspace{2mm}\hrule\vspace{2mm}
}

\newcommand{\job}[4]{%
    \textbf{\large{#1/}}\textcolor{Plum}{#2}\hfill\textbf{\underline{#3}}
}

\newcommand{\edu}[3]{%
    \noindent\textbf{\large{#1/}}\textcolor{Plum}{#2}\hfill\textbf{\underline{#3}}
}

\newcommand{\pubfirst}[3]{%
    \vspace{1mm}
    \noindent
    \begin{minipage}[t]{24mm}
        \textbf{\large #1}
    \end{minipage}%
    \hspace{3mm}%
    \begin{minipage}[t]{\dimexpr\linewidth-30mm\relax}
        {\large\textbf{#2}}\par
        \vspace{0.8mm}
        #3
    \end{minipage}
    \par\vspace{3mm}
}
\newcommand{\pubnext}[2]{%
    \noindent
    \hspace*{27mm}%
    \begin{minipage}[t]{\dimexpr\linewidth-33mm\relax}
        {\large\textbf{#1}}\par
        \vspace{0.8mm}
        #2
    \end{minipage}
    \par\vspace{3mm}
}


\setromanfont{LeagueSpartan}[
    Path=./Font/,
    Extension = .ttf,
    UprightFont=*-Regular,
    BoldFont=*-SemiBold
    ]

\author{Vladimir Baikalov}

\setlength\columnsep{10mm}
\setlength\parindent{0pt}
\pagenumbering{gobble}

\begin{document}

\noindent
\begin{minipage}[c]{0.65\linewidth}
\Huge\textbf{Vladimir Baikalov}

\small{Research Engineer}

\small{Generative Retrieval, Continual Learning, Foundation Recommender Systems}

\end{minipage}
\begin{minipage}[c]{0.35\linewidth}
\raggedleft

\href{mailto:nonameuntitled159@gmail.com}{nonameuntitled159@gmail.com}\\

\href{https://www.linkedin.com/in/noname-untitled}{linkedin.com/in/noname-untitled}\\

\href{http://www.github.com/NonameUntitled}{github.com/NonameUntitled}

\href{https://scholar.google.com/citations?user=VNlOjZEAAAAJ}{\underline{Google Scholar Link}}

\end{minipage}

\vspace{2mm}

\cvsection{Profile}

Research engineer building foundation-scale generative recommender and search/retrieval systems that integrate multimodal signals, adapt under non-stationarity, and optimize long-horizon user outcomes.

\vspace{2mm}

5+ years of production ML experience across Google/YouTube, Yandex, AI VK, and Huawei, taking models from research prototypes to large-scale deployment. PhD candidate with publications at WWW, RecSys, SIGIR, KDD.

\vspace{2mm}

\textbf{Skills}: Python, C++, PyTorch, JAX, TensorFlow, Distributed ML, GPU Acceleration, Large-Scale Data Pipelines, Docker, Kubernetes.

\vspace{2mm}

\textbf{Research Interests}: Generative Retrieval, Foundation RecSys, World Models, Bayesian Statistics, Continual Learning, Offline RL, Long-Term Objectives, Multimodal Representation Learning.

\vspace{1mm}

\cvsection{Work experience}

\job{Lead RecSys Scientist}{AI VK}{Dec 2025 -- Present}

\begin{itemize}[itemsep=0.0mm, topsep=1mm]

    \item Leading research on foundation recommender systems for a large video platform, with a focus on generative retrieval, continual semantic updates, and reward-aware long-horizon optimization.
    
    \item Co-authored two research works on adaptive generative recommendation: mitigating Semantic ID staleness in generative retrieval and off-policy optimization for long-horizon objectives in large action spaces.

    \item Developing generative retrieval models for video service and preparing an online A/B test in collaboration with the product team.

    \item Developing a unified recommender foundation model that combines generative modeling, LLM $\times$ RecSys, and reward-aware training.
    
\end{itemize}

\vspace{2mm}

\job{Lead, R\&D RecSys Team}{Yandex Technologies}{Jun 2024 -- Nov 2025}

\begin{itemize}[itemsep=0.0mm, topsep=1mm]
    
    \item Led a three-person R\&D team building generative end-to-end recommender models on semantic IDs as an alternative to classic multi-stage retrieval/ranking pipelines.

    \item Co-developed Argus, increasing Total Listening Time at Yandex.Music by +2.26\% It was later adapted to other Yandex products (e.g., e-commerce +1.7\% GMV, food-tech services +2.1\% GMV).

    \item Full-stack AI engineer: Owned entire end-to-end ML pipeline from data collection to inference. Built data collection, training, evaluation, and deployment for retrieval \& ranking models.
    
\end{itemize}

\vspace{2mm}

\job{Machine Learning Engineer}{Google, YouTube}{Jan 2023 -- Jun 2024}

\begin{itemize}[itemsep=0.0mm, topsep=1mm]
    \item Developed and launched a detection model for \textbf{YouTube Shorts}, improving recall by 4\% without precision degradation and enabling more frequent fine-tuning of content-quality pipelines.

    \item Designed an end-to-end serving system for this classifier in collaboration with the \textbf{YouTube Gaming} team, taking the model from research prototype to production deployment.

    \item Worked on production deployment of LLM-based systems, including improvements to fine-tuning and distillation pipelines for large models.

    \item Built ML systems under \textbf{Google engineering practices}, including design/code reviews, production reliability constraints, and cross-team collaboration on large codebases.
    
\end{itemize}

\vspace{2mm}

\job{Deep Learning Engineer}{Yandex Technologies}{Aug 2021 -- Dec 2022}

\begin{itemize}[itemsep=0.0mm, topsep=1mm]
    \item Ported a research PyTorch model into a production C++ ML framework for Ads, increasing GMV by up to 1.5\% and CTR by up to 5\%.

    \item Built a distributed parallel data-processing pipeline, improving GPU utilization by up to 2.5 times.
    
\end{itemize}

\vspace{2mm}

\job{Machine Learning Engineer}{Huawei R\&D Dept.}{Mar 2020 -- Jun 2021}

\vspace{1mm}

\cvsection{Education}

\edu{PhD Student}{ITMO University}{Sep 2024 -- Present}

{\leftskip=3mm
    
    Research: Methods for training, scaling, and updating representations in generative recommender systems. 

    \vspace{2mm}
}



\edu{Master Degree}{Skoltech, DS major}{Sep 2021 -- Jul 2024}

{\leftskip=3mm

    Honors Thesis: Enhancing recommender systems performance with multiple data representations.

    \vspace{2mm}
}

\edu{Bachelor Degree}{ITMO University, CS major}{Sep 2017 -- Aug 2021}

{\leftskip=3mm

    Thesis: Application of Multi-Agent Routing Algorithms Based on Policy Optimization.

    \vspace{2mm}
}

\cvsection{Publications}

\vspace{1mm}

\pubfirst{KDD'26}
{\href{https://arxiv.org/abs/2507.15994}{Scaling Recommender Transformers to One Billion Parameters}}
{Kirill Khrylchenko, Artem Matveev, Sergei Makeev, \underline{Vladimir Baikalov}}

\pubnext
{\href{https://arxiv.org/abs/2607.02818}{Long-Term Optimization for Large-Scale Generative Retrieval with Off-Policy REINFORCE}}
{Artem Matveev, Sergei Makeev, Aleksei Krasilnikov, \underline{Vladimir Baikalov}, Sergei Liamaev, Kirill Khrylchenko}

\pubfirst{SIGIR'26}
{\href{https://arxiv.org/abs/2604.13273}{Mitigating Collaborative Semantic ID Staleness in Generative Retrieval}}
{\underline{Vladimir Baikalov}, Iskander Bagautdinov, Sergey Muravyov}

\pubfirst{RecSys'25}
{\href{https://dl.acm.org/doi/abs/10.1145/3705328.3748033}{Correcting the LogQ Correction: Revisiting Sampled Softmax for Large-Scale Retrieval}}
{Kirill Khrylchenko, \underline{Vladimir Baikalov}, Sergei Makeev, Artem Matveev, Sergei Liamaev}

\pubnext
{\href{https://dl.acm.org/doi/abs/10.1145/3705328.3748163}{Yambda-5B: A Large-Scale Multi-Modal Dataset for Ranking and Retrieval}}
{Alexander Ploshkin, Vladislav Tytskiy, Alexey Pismenny, \underline{Vladimir Baikalov}, Evgeny Taychinov, Artem Permiakov, Daniil Burlakov, Eugene Krofto}

\pubnext
{\href{https://dl.acm.org/doi/full/10.1145/3758126.3758131}{Blending Sequential Embeddings, Graphs, and Engineered Features: 4th Place Solution in RecSys Challenge 2025}}
{Sergei Makeev, Alexandr Andreev, \underline{Vladimir Baikalov}, Vladislav Tytskiy, Aleksei Krasilnikov, Kirill Khrylchenko}

\pubfirst{WWW'24}
{\href{https://arxiv.org/abs/2403.00895}{End-to-End Graph-Sequential Representation Learning for Accurate Recommendations}}
{\underline{Vladimir Baikalov}, Evgeny Frolov}

\pubfirst{Elsevier'21}
{\href{https://aaltodoc.aalto.fi/handle/123456789/111642}{Multi-Agent Deep Reinforcement Learning-Based Algorithm for Fast Generalization on Routing Problems}}
{Ibraheem Barbahan, \underline{Vladimir Baikalov}, Valeriy Vyatkin, Andrey Filchenkov}

\vspace{1mm}

\cvsection{Teaching Experience}

\textbf{\large{Teaching Assistant/}}\textcolor{Plum}{Yandex School of Data Analysis (YSDA)}\hfill\textbf{\underline{2024 -- 2026}}

Recommender System Course.
\textcolor{Blue}{\underline{\href{https://github.com/yandexdataschool/recsys_course}{GitHub link (140+ stars)}}}

\vspace{2mm}

\textbf{\large{Teaching Assistant/}}\textcolor{Plum}{Higher School of Economics}\hfill\textbf{\underline{2025 -- 2026}}

Recommender System Course. \textcolor{Blue}{
\underline{
\href{https://github.com/KhrylchenkoKirill/DeepRecSys}{GitHub link (90+ stars)}
}
}

\vspace{2mm}

\textbf{\large{Lecturer/}}\textcolor{Plum}{ITMO University}\hfill\textbf{\underline{2024 -- 2026}}

Taught Programming Massively Parallel Processors.

Topics include: GPU architecture, CUDA, Triton, Kernels, Inference Acceleration, Large-Scale Deep Learning.

\vspace{1mm}

\end{document}
